\subsection{مبانی نظری و فرمول‌بندی ریاضی}
در روش‌های آگاه از فعال‌سازی، ماتریس وزن با جذر ماتریس کوواریانس فعال‌سازی‌ها مقیاس‌گذاری می‌شود: $\mathbf{W}_{\text{weighted}} = \mathbf{W} \mathbf{C}_X^{1/2}$. در فشرده‌ساز \lr{DRONE}، این فرض مبنایی قرار می‌گیرد که فعال‌سازی‌های ورودی سفید و ناهمبسته بوده و دارای کوواریانس همانی هستند:
\begin{equation}
\mathbf{C}_X = \mathbf{I} \implies \mathbf{W}_{\text{weighted}} = \mathbf{W}
\end{equation}

بدین ترتیب، تجزیه مقادیر منفرد مستقیماً بر روی ماتریس وزن اولیه $\mathbf{W} \in \mathbb{R}^{d_{\text{out}} \times d_{\text{in}}}$ اجرا می‌گردد:
\begin{equation}
\mathbf{W} = \mathbf{U} \mathbf{\Sigma} \mathbf{V}^\top
\end{equation}

رتبه فشرده‌سازی $r$ با استفاده از تابع پویای تخصیص رتبه با اعمال قید کاهش قطعی پارامترها تعیین می‌شود:
\begin{equation}
r = \min\left(r_{\text{target}}, d_{\text{in}}, d_{\text{out}}, \left\lfloor 0.75 \frac{d_{\text{in}} \cdot d_{\text{out}}}{d_{\text{in}} + d_{\text{out}}} \right\rfloor \right)
\end{equation}

مقادیر منفرد $\mathbf{\Sigma}_r$ با جذر متقارن به فاکتورهای چپ و راست تخصیص داده می‌شوند:
\begin{equation}
\mathbf{U}_r = \mathbf{U}_{:, :r} \, \mathbf{\Sigma}_r^{1/2}, \quad \mathbf{V}_r = \mathbf{\Sigma}_r^{1/2} \, \mathbf{V}_{:, :r}^\top
\end{equation}

ماتریس بازسازی‌شده $\mathbf{W}_{\text{recon}} = \mathbf{U}_r \mathbf{V}_r$ مورد ارزیابی خطای فروبنیوس قرار می‌گیرد:
\begin{equation}
\epsilon_{\text{fro}} = \frac{\|\mathbf{W} - \mathbf{W}_{\text{recon}}\|_F}{\max(\|\mathbf{W}\|_F, 10^{-12})}
\end{equation}
در صورتی که $\epsilon_{\text{fro}} > \epsilon_{\text{max}}$ (با آستانه مجاز $\epsilon_{\text{max}} = 0.50$) باشد، فشرده‌سازی لایه لغو شده و ساختار اولیه حفظ می‌گردد. در غیر این صورت، لایه خطی با ماژول دو مرحله‌ای جایگزین می‌شود:
\begin{equation}
\mathbf{y} = (\mathbf{x} \mathbf{V}_r^\top) \mathbf{U}_r^\top + \mathbf{b}
\end{equation}
